% titel: Blatt 3, A3: Grenzwerte von Folgen

% Teil I: Werkzeuge, kein Lösungsweg.

Grenzwerte von Folgen für $n\to\infty$ bestimmen: rationale Folgen,
Wurzelbrüche, Exponential-/rationale Kombinationen und Fakultäten.

\begin{werkzeugkasten}[Führende Terme]
Bei Quotienten von Polynomen dominiert die höchste Potenz:
\begin{equation}
  \frac{3n^3-2n^2+1}{n^2-2n^3-3}\ \sim\ \frac{3n^3}{-2n^3}
  \to -\frac32 .
\end{equation}
Bei $2^n,3^n,\dots$ dominiert die größte Basis; $n!\ $ wächst schneller als
jede feste Potenz:
\begin{equation}
  \frac{n^2}{(n+1)!}\to0 .
\end{equation}
\end{werkzeugkasten}

\begin{werkzeugkasten}[Wurzelbrüche]
Beim Grenzwert mit Wurzeln im Nenner durch die führende Potenz von $n$
kürzen (z. B. durch $\sqrt n$); dabei $\sqrt{n+1}\approx\sqrt n$ nutzen,
das Restglied $\to0$.
\end{werkzeugkasten}

\begin{achtung}[Darauf achten]
\begin{itemize}
  \item Kürzen durch die \emph{größte} auftretende Potenz von $n$; die
        kleineren Terme fallen dann weg.
  \item Bei $(-1)^n$ gemischten Termen bleibt das Vorzeichen des
        führenden Glieds entscheidend; $(-1)^n/3^n\to0$.
  \item $3^{n+1}=3\cdot3^n$ -- im Nenner $4\cdot3^{n+1}$ die Basis $3$ mit
        $3^n$ vergleichen, nicht $2^n$.
  \item $d_n=n^2/(n+1)!$: Fakultät dominiert, Grenzwert $0$.
\end{itemize}
\end{achtung}
